%! Author = julian
%! Date = 19.04.26


\section{Model Comparison} \label{sec:model-comparison}
If different models both classify a portion of data as preictal, even though no seizure occurs afterward, this is interpreted as evidence for a proictal state, since preictal features were apparently present in the data.
To quantify the tendency of false predictions to co-occur across models, we calculated the \gls{cope} \parencite{muller_coherent_2022} between the ensemble and \gls{cnn} scores for each patient.
Its definition follows.

Let $\vec{p}_i$ and $\vec{p}_j$ be the model's outputs for models $i$ and $j$, respectively, and $\vec{L}$ the correct label vector.
The prediction error vector of model $i$ is defined as $\vec{e}_i=\lvert\vec{p}_i-\vec{L}\rvert$ (analogously for model $j$).
The weight $w_k$ for each sample $k$, the weighted mean $\langle z\rangle_w$, and weighted standard deviation $\sigma_w(z)$ are defined as:
\begin{equation} \label{eq:cope-helpers}
    \begin{aligned}
        w_k &= \max(e_{i,k},e_{j,k}),\\
        \langle z\rangle_w &= \frac{\sum_k w_k z_k}{\sum_k w_k},\\
        \sigma_w(z) &= \sqrt{\langle z^2\rangle_w-\langle z\rangle_w^2}.
    \end{aligned}
\end{equation}
The \gls{cope} is then
\begin{equation} \label{eq:cope}
    c_w=\frac{\left\langle\left(\vec{e}_i-\langle \vec{e}_i\rangle_w\right)\left(\vec{e}_j-\langle \vec{e}_j\rangle_w\right)\right\rangle_w}
    {\sigma_w(\vec{e}_i)\,\sigma_w(\vec{e}_j)}.
\end{equation}

This measure emphasizes coherence in false predictions by assigning higher weight to samples where at least one model has a large prediction error.
Consequently, samples with low errors contribute little, whereas samples with high errors dominate the estimate.